\chapter{Diskussion}
\label{ch:diskussion}

\todo[inline]{Vorläufige Fassung auf dem Stand vom 19.11. Nach den Versuchen mit dem neuen
Datensatz überarbeiten.}

% ==============================================================================
\section{Bewertung der Designentscheidungen}
\label{sec:bewertung}

\Cref{tab:bewertung} stellt die wichtigsten Entscheidungen ihrer Wirkung in den Versuchen
gegenüber.

\begin{table}[htbp]
  \centering
  \caption{Designentscheidungen und ihre Bewertung nach den Versuchen}
  \label{tab:bewertung}
  \small
  \begin{tabularx}{\textwidth}{@{}>{\raggedright\arraybackslash}p{38mm} X@{}}
    \toprule
    Entscheidung & Bewertung \\
    \midrule
    Ports und Adapter, hardwarefrei entwickelt & bewährt: Die Kette lief vor dem ersten
      Anlagentag vollständig; an der Anlage blieben nur Adapter- und Physikthemen. \\
    Planer statt Teleoperation & bewährt: gleichmäßiges Tempo und über 700 Demonstrationen
      mit wenig Handarbeit. Nachteil: Alle Bahnen haben dieselbe Struktur; was der Planer nie
      zeigt (z.\,B.\ Drehen unter \code{ANFLUG}), kann die Policy nicht. \\
    Rauschen mit Trichter & in der Vorstudie deutlich wirksam. An der Anlage zeigte sich die
      Lücke im Endanflug, die der gezielte Versatz schließen soll. \\
    Absolute Gelenkwinkel als Aktion & bewährt: direkte Abbildung auf \code{servo_j}, keine
      Drift durch aufsummierte Änderungen. \\
    Direkte Vorhersage der Aktionsfolge & notwendig: Mit Rauschvorhersage waren zehn
      Entrauschungsschritte zu ungenau. \\
    \code{servo_j} mit Override \num{0.5} & bewährt: ruhig wie ein Herstellerprogramm und
      bahntreu. \\
    Keine Rektifizierung & bestätigt: Die Policy reagiert kaum auf die Kameralage. \\
    Eigene Sicherheitsschicht & notwendig: In der \acs{vm} wäre die Werkzeugspitze in gut der
      Hälfte der Fahrten unter den Tisch gegangen. Die Stufen sind jeweils aus einem
      beobachteten Fehlerfall entstanden. \\
    Policy nur für den Griff & vielversprechend: weniger Abbrüche und Tischfahrten, gleiche
      Genauigkeit. \\
    \bottomrule
  \end{tabularx}
\end{table}

% ==============================================================================
\section{Grenzen}
\label{sec:grenzen}

\paragraph{Aussagekraft der Tests ohne Anlage.} Weder der Offline-Test noch der \ac{vm}-Test
koppeln die Bilder an die Bewegung zurück. Sie eignen sich zum Vergleich von Modellen und zum
Finden grober Fehler, sagen aber die Genauigkeit an der Anlage nur ungefähr und den Drehfehler
gar nicht vorher. Der einzige vollständige Test ist die Anlage, an der aus Zeit- und
Sicherheitsgründen nur wenige Fahrten möglich waren.

\paragraph{Fehlerfortpflanzung bleibt.} Offline ist die Policy auf wenige Millimeter genau, an
der Anlage auf knapp einen Zentimeter. Der Unterschied entsteht in Zuständen, die sie selbst
herbeiführt. Das ist genau die Hürde aus \cref{sec:covariate-shift}; das eingespeiste
Rauschen verringert sie, beseitigt sie aber nicht dort, wo es gedämpft ist.

\paragraph{Ungeklärte Verschlechterung.} Warum die Daten vom 11.11.\ die Genauigkeit
verschlechterten, ist offen. Der Befund zeigt, wie empfindlich das Ergebnis auf schwer
greifbare Unterschiede der Aufnahmebedingungen reagiert.

\paragraph{Umfang.} Alle Versuche verwenden ein Bauteil auf einem ebenen Tisch, ohne Kiste und
ohne Störungen während der Fahrt. Aussagen über andere Bauteile oder Szenen sind daraus nicht
ableitbar.
